\documentclass[12pt,a4paper,oneside]{book}

% ---------- PREAMBLE ----------
\usepackage[margin=25mm]{geometry}
\usepackage[T1]{fontenc}
\usepackage{graphicx}
\usepackage{amsmath,amssymb}
\usepackage{array,tabularx,multirow}
\usepackage{enumitem}
\usepackage{multicol}
\usepackage{setspace}
\usepackage{xcolor}
\usepackage{makeidx}
\usepackage[colorlinks=true,linkcolor=blue,urlcolor=blue,citecolor=blue]{hyperref}

\makeindex
\setlength{\parskip}{5pt}

\title{\Huge Artificial Intelligence\\[6pt]
       \Large Concepts, Techniques and Applications}
\author{Your Name}
\date{\today}

\begin{document}

% ---------- FRONT MATTER ----------
\frontmatter
\maketitle

\chapter*{Preface}
This sample book is designed to demonstrate the LaTeX features in the
course syllabus. Replace this text with your own preface.

\tableofcontents
\listoffigures
\listoftables

% ---------- MAIN MATTER ----------
\mainmatter

\chapter{Introduction to Artificial Intelligence}
Artificial Intelligence (AI) is the study of computer systems that
perform tasks normally associated with human intelligence.
The term AI is indexed here\index{Artificial Intelligence}.

\section{Basic idea}
A computer system receives input, processes it, and produces output.
Figure~\ref{fig:workflow} illustrates this simple idea.

\begin{figure}[htbp]
\centering
\includegraphics[width=0.8\linewidth]{workflow.png}
\caption{Input-process-output workflow}\label{fig:workflow}
\end{figure}

\section{Formatting}
This is \textbf{bold}, this is \textit{italic}, and this is
\underline{underlined}. A footnote can explain a term.\footnote{A
footnote appears at the bottom of the page.}

\chapter{Lists and Tables}
\section{Types of lists}
\begin{itemize}
  \item Supervised learning
  \item Unsupervised learning
  \item Reinforcement learning
\end{itemize}

\begin{enumerate}[label=\Roman*.]
  \item Data collection
  \item Model training
  \item Model evaluation
\end{enumerate}

\begin{description}
  \item[Algorithm] A sequence of steps to solve a problem.
  \item[Dataset] A collection of examples.
\end{description}

\section{A sample table}
Table~\ref{tab:areas} summarises example applications.
\begin{table}[htbp]
\centering
\caption{AI application areas}\label{tab:areas}
\begin{tabular}{|l|l|}
\hline
\textbf{Area} & \textbf{Example}\\ \hline
Healthcare & Medical image analysis\\ \hline
Education & Personalised learning\\ \hline
Transport & Route planning\\ \hline
\end{tabular}
\end{table}

\chapter{Mathematics and Cross-References}
A displayed equation:
\begin{equation}\label{eq:neuron}
z=\sum_{i=1}^{n}w_i x_i+b
\end{equation}
Equation~\eqref{eq:neuron} is a weighted sum plus a bias term.

An aligned example:
\begin{align}
(a+b)^2 &= (a+b)(a+b)\\
        &= a^2+2ab+b^2
\end{align}

See Figure~\ref{fig:workflow} on page~\pageref{fig:workflow}.

\chapter{Large Documents and Child Files}
A large book can be divided into separate source files. For example,
\texttt{chapter-extra.tex} can be inserted with
\verb|\chapter{Extra Chapter Inserted from a Separate File}
This chapter demonstrates the use of \verb|\chapter{Extra Chapter Inserted from a Separate File}
This chapter demonstrates the use of \verb|\chapter{Extra Chapter Inserted from a Separate File}
This chapter demonstrates the use of \verb|\input{chapter-extra}|.
The child file contains chapter content only; it does not contain
\verb|\documentclass| or \verb|\begin{document}|.
|.
The child file contains chapter content only; it does not contain
\verb|\documentclass| or \verb|\begin{document}|.
|.
The child file contains chapter content only; it does not contain
\verb|\documentclass| or \verb|\begin{document}|.
|. The main file retains the preamble and
document structure.

\appendix
\chapter{Quick Reference}
\begin{tabularx}{\linewidth}{|l|X|}
\hline
Command & Purpose\\ \hline
\verb|\chapter{}| & Starts a chapter\\
\verb|\section{}| & Starts a section\\
\verb|\label{}| & Defines a reference key\\
\verb|\ref{}| & Prints the referenced number\\
\hline
\end{tabularx}

% ---------- BACK MATTER ----------
\backmatter
\chapter*{Bibliography}
This demonstration uses a manually written back-matter heading.
For an automatically generated BibTeX bibliography, use:
\verb|\bibliographystyle{plain}| and
\verb|\bibliography{references}| before \verb|\printindex|.

\printindex
\end{document}
